% STAGED 2026-09-25. Replaces the B9 paragraph and placeholder figure only.
% Upload the staged steeleknee.pdf before applying this block.

Steele's biomimetic knee joint supplies the physical knee used in this dissertation \citep{steele_development_2017, steele_biomimetic_2018, steele_biomimetic_2018-1}. It is a one-degree-of-freedom crossed four-bar linkage that reproduces the posterior migration of the human knee's instantaneous center of rotation (ICR) during flexion. The linkage geometry therefore produces rolling-and-sliding motion without cams or gear profiles. Because the ICR migrates, the moment arms of muscles crossing the knee vary with flexion angle, and a torque prediction must update both the force direction and moment arm throughout the range of motion. The biomimetic-knee experiments in Sections~\ref{sec:robot_model}, \ref{sec:results_torque}, and \ref{sec:disc_error} test the corrected torque model against this geometry.

\begin{figure}[htbp]
\centering
\includegraphics[width=0.95\textwidth]{steeleknee.pdf}
% ALT TEXT: Two panels reproduced from Steele et al. Figure 6. Panel (a) is a
% sagittal cutaway of the biomimetic knee at full extension, showing the crossed
% four-bar linkage. Panel (b) is an exploded view of the knee's mechanical stop.
\caption[Steele biomimetic knee mechanism]{Biomimetic knee mechanism. \textbf{(a)} Sagittal cutaway at full extension showing the crossed four-bar linkage. \textbf{(b)} Exploded view of the mechanical stop. Reproduced from \citet[Fig.~6]{steele_experimental_2023} under CC BY 4.0.}
\label{fig:steele_knee}
\end{figure}
